% Chapter Template

\chapter{Ensayos y resultados} % Main chapter title

\label{Chapter4} % Change X to a consecutive number; for referencing this chapter elsewhere, use \ref{ChapterX}
Todos los capítulos deben comenzar con un breve párrafo introductorio que indique cuál es el contenido que se encontrará al leerlo.  La redacción sobre el contenido de la memoria debe hacerse en presente y todo lo referido al proyecto en pasado, siempre de modo impersonal.

%----------------------------------------------------------------------------------------
%	SECTION 1
%----------------------------------------------------------------------------------------

\section{Consideraciones generales}

xx

%----------------------------------------------------------------------------------------
\section{FEINN}

xx

\subsection{Caso de estudio 1: Placa con concentrador de tensiones}

xx

\subsection{Caso de estudio 2: Matriz blanda con inclusión}

xx

%----------------------------------------------------------------------------------------
\section{Modelo subrogado basado en FNO}

xx

\subsection{Configuración del caso de estudio}

xx

\subsection{Generación de datos sintéticos}

xx

\subsection{Optimización de hiperparámetros}

xx

\subsection{Entrenamiento del modelo}

xx

\subsection{Evaluación de métricas}

xx

%----------------------------------------------------------------------------------------
\section{Acoplamiento FEINN-FNO para análisis multiescala}

xx

\subsection{Análisis de resultados}

xx
